\section{Theoretical Analyses}
\label{sec:theory-main}

We now present theoretical analyses of the LDM. This section is self-contained; readers primarily interested in the empirical results may safely skip it. Let \(\X\) denote the design space, and let \(x^\star \in \argmax_{x \in \X} R(x)\) be a global maximiser of the unknown objective function. 
Note that performance is measured with respect to the true objective \(R\), rather than the uncertainty-aware acquisition function. After \(T\) trials, the simple regret (no separate symbol beyond the per-round regret) is defined as
\begin{equation}
\regret_T = R(x^\star) - \max_{1 \le t \le T} R(x_t).
\end{equation}

At round \(t\), let \(p_\theta(\cdot\mid\C_t)\) denote the LLM reservoir distribution after applying the actual decoding restrictions of §method (top-\(k\) or top-\(p\) truncation, etc.). The corresponding temperature-adjusted proposal is
\[
p_{\theta,\alpha}(x\mid\C_t) \;:=\;
\frac{p_\theta(x\mid\C_t)^\alpha}
{\sum_{u \in \X} p_\theta(u\mid\C_t)^\alpha}\,,
\]
and its support (the LLM-controlled dynamic search domain of §\ref{sec:algo}) is
\[
A_t \;:=\; \operatorname{supp}\, p_{\theta,\alpha}(\cdot\mid\C_t) \;=\; \{x\in\X : p_{\theta,\alpha}(x\mid\C_t) > 0\}.
\]

The LDM then samples the next candidate \(x_t \sim \pi_t\), where
\begin{equation}
\label{eq:theory-pi}
\pi_t(x) =
\frac{\exp(\eta\,\acq_t(x))\,p_{\theta,\alpha}(x\mid\C_t)}
{\sum_{u \in A_t} \exp(\eta\,\acq_t(u))\,p_{\theta,\alpha}(u\mid\C_t)}.
\end{equation}
Here, $\acq_t(x) = \mu_t(x) + \sqrt{\beta_t}\sigma_t(x)$ is the UCB acquisition function (matching the LDM convention of §method).

To relate this performance criterion to the per-round behaviour of the algorithm, define the instantaneous regret at round \(t\) as $\regret_t = R(x^\star) - R(x_t)$. Here, \(t\) indexes evaluations under the true objective \(R\), rather than internal LLM interaction steps. Equivalently, the simple regret after \(T\) trials is $\regret_T = \min_{1 \le t \le T} \regret_t$. In the analysis below, we first study the cumulative behaviour of \(\{\regret_t\}_{t=1}^T\), from which a simple regret guarantee follows via the standard relation $\regret_T \le \frac{1}{T}\sum_{t=1}^T \regret_t$.

\noindent\textbf{Instantaneous Regret Decomposition as Discovery Gap and Optimisation Regret.}
For each round \(t\), define the best objective value attainable within the current reservoir support as $R_{t,A}^\star = \max_{x \in A_t} R(x)$. Then, the instantaneous regret admits the exact decomposition
\begin{align}
\regret_t
&= R(x^\star) - R(x_t) \\
&=
\underbrace{R(x^\star) - R_{t,A}^\star}_{\regret_t^{\rm disc}}
+
\underbrace{R_{t,A}^\star - R(x_t)}_{\regret_t^{\rm opt}} .
\end{align}
The first term, \(\regret_t^{\rm disc}\), is the \emph{discovery gap}: it measures the loss incurred because the current reservoir support \(A_t\) may not contain designs sufficiently close to the global optimum. This term is zero whenever \(x^\star \in A_t\). The second term, \(\regret_t^{\rm opt}\), is the \emph{optimisation regret} within the currently discovered region: it measures how far the sampled point \(x_t\) is from the best candidate available in \(A_t\).

This decomposition separates the two roles of the LDM. The reservoir distribution controls whether high-quality regions are discovered at all, while the acquisition tilt controls how effectively the algorithm selects promising candidates among those currently reachable.

Below, we will give the necessary assumptions and definitions to support the regret bound analysis. 
\begin{assumption}[UCB calibration]
\label{ass:ucb-calibration}
For a non-decreasing sequence $(\beta_t)_{t\ge1}$, with probability at least $1-\delta_{\rm gp}$, simultaneously
for all $t\le T$ and $x\in\X$,
\[
|R(x)-\mu_t(x)|
\le
\sqrt{\beta_t}\,\sigma_t(x).
\]
This event is implied by the standard kernel bandit conditions: $R(x)=m(x)+g(\phi(x))$ with $g\in\mathcal H_k$, bounded RKHS norm, bounded kernel diagonal, and conditionally sub-Gaussian observation noise \citep{srinivas2010gaussian,chowdhury2017kernelized}.
\end{assumption}

\begin{assumption}[Reservoir coverage and local regularity]
\label{ass:coverage}
Let $\mathrm d_\X$ be a distance on the design space and define the reservoir
coverage radius $r_t^{\rm cov}=\inf_{x\in A_t}\mathrm d_\X(x,x^\star)$. The objective is one-sided H\"older around the optimum: there exist $L>0$ and $\chi\in(0,1]$ such that, for all
$x\in\X$,
\[
R(x^\star)-R(x)\le L\,\mathrm d_\X(x,x^\star)^\chi.
\]
\end{assumption}
Let $a_{t,A}^\star=\sup_{x\in A_t}\acq_t(x)$ be the maximum value of the acquisition function at $t$-th round. Then, for any tolerance $\zeta_t\ge0$, let $U_t(\zeta_t)=
\{x\in A_t:\acq_t(x)\ge a_{t,A}^\star-\zeta_t\}$ be the set of candidates in the current reservoir support $A_t$ whose acquisition value is within $\zeta_t$ of the best acquisition value. We denote the reservoir probability mass assigned to this near-best set by
\begin{equation}
\label{eq:coverage}
\kappa_t(\zeta_t)=p_{\theta,\alpha}\bigl(U_t(\zeta_t)\mid\C_t\bigr).
\end{equation}
The bound below depends on the size of $\kappa_t(\zeta_t)$: a small near-UCB mass means that the LLM reservoir makes high-acquisition designs hard to sample even after exponential tilting.

\begin{equation}
\label{eq:adaptive}
\kappa_t(\zeta)\;\ge\;\kappa_0(\zeta)\;+\;\rho_\zeta\cdot N_t^{\rm good},
\end{equation}
where $N_t^{\rm good}:=|\{(x_i,r_i)\in\C_t : r_i \ge r^{\rm thresh}\}|$ counts high-reward exemplars in the
context. Equation~\eqref{eq:adaptive} formalises the adaptive condition: the near-acquisition reservoir mass grows monotonically as good exemplars
accumulate.

\begin{lemma}[Gibbs near-UCB sampling]
\label{lem:gibbs-near-ucb}
Conditioned on the history \(\C_t\), fix any \(\zeta_t\ge 0\) such that the near-UCB set has reservoir mass \(\kappa_t(\zeta_t)\). Then, for any \(\rho_t\in(0,1)\), we have
\begin{equation}
 \mathbb P\!\left(
a_{t,A}^\star-\acq_t(x_t)
\le
\zeta_t+\frac{1}{\eta}\log\frac{1}{\rho_t\kappa_t(\zeta_t)}
\,\middle|\, \C_t
\right)
\ge
1-\rho_t .
\end{equation}

\end{lemma}

Let
\[
\gamma_T=
\sup_{B\subset\X,\ |B|\le T}
\frac12\log\det(I+\lambda^{-1}K_B),
\qquad
C_\lambda=\frac{2}{\log(1+\lambda^{-1})},
\]
where $K_B$ is the kernel matrix on $B$.

\begin{theorem}[Simple regret of large discovery model]
\label{thm:regret}
Suppose Assumptions~\ref{ass:ucb-calibration} and~\ref{ass:coverage} hold. Fix $\delta_{\rm gp},\delta_{\rm samp}\in(0,1)$ and choose
$\rho_t>0$ such that $\sum_{t=1}^T\rho_t\le\delta_{\rm samp}$. Then, for any $\zeta_t\ge0$ with
$\kappa_t(\zeta_t)>0$, with probability at least $1-\delta_{\rm gp}-\delta_{\rm samp}$,
\[
  \boxed{
  \frac{1}{T}\sum_{t=1}^T \regret_t
  \le
  \underbrace{\frac{L}{T}\sum_{t=1}^T(r_t^{\rm cov})^\chi}_{\text{discovery gap}}
  +
  \underbrace{2\sqrt{\frac{C_\lambda\beta_T\gamma_T}{T}}}_{\text{GP-UCB optimisation}}
  +
  \underbrace{
  \frac1T\sum_{t=1}^T
  \left[
  \zeta_t+
  \frac1\eta\log\frac1{\rho_t\kappa_t(\zeta_t)}
  \right]
  }_{\text{LDM sampling shortfall}} .
  }
\]
\end{theorem}

The first term is the gap of not yet having generated candidates near the true optimum. The second is the usual GP-UCB information-gain term within the covered region. The third is the cost of sampling from the tilted LDM distribution rather than exactly maximising UCB over $A_t$; it decreases when the acquisition tilt $\eta$ is sharper or when the near-UCB reservoir mass $\kappa_t(\zeta_t)$ is larger.

The LLM affects the regret bound through the reservoir support \(A_t\) and its distribution \(p_{\theta,\alpha}(\cdot\mid\C_t)\). First, an informative LLM can propose candidates closer to the global optimum, thereby reducing the coverage radius \(r_t^{\rm cov}\) and hence the discovery-gap term. Second, it can assign substantial probability mass to candidates with high acquisition values, increasing \(\kappa_t(\zeta_t)\) and reducing the LDM sampling-shortfall term. Thus, the LLM provides a structured, history-dependent proposal distribution that can focus evaluations on semantically plausible and promising regions of the design space, rather than searching the entire space uniformly.
